\documentclass[12pt,a4paper]{article}
\usepackage[utf8]{inputenc}
\usepackage[T1]{fontenc}
\usepackage{lmodern}
\usepackage{geometry}
\usepackage{longtable}
\usepackage{hyperref}
\usepackage{enumitem}
\geometry{margin=2.5cm}

\title{Revisão Arquitetural e Consolidação da RDC}
\author{Projeto RDC}
\date{\today}

\begin{document}
	\maketitle
	
	\section{Análise Arquitetural}
	
	Excelente análise — você fez exatamente o papel de \textbf{revisor arquitetural}, não apenas de autor.
	Os ajustes apontados foram \textbf{reais, objetivos e necessários} e já foram aplicados ao documento oficial.
	
	Este texto resume as alterações realizadas e alinha o \textbf{próximo passo correto}, mantendo a disciplina arquitetural estabelecida.
	
	\section{Ajustes Aplicados com Sucesso}
	
	\subsection{Introdução Reestruturada (Consistência Interna)}
	
	A duplicação foi eliminada e a abertura do documento passou a seguir um \textbf{padrão arquitetural clássico e limpo}:
	
	\begin{itemize}[leftmargin=2em]
		\item \textbf{1. Introdução} — Define o que é a RDC, sem redundância.
		\item \textbf{1.1 Propósito} — Explicita o papel da SPEC como \emph{fonte de verdade} e artefato governado.
		\item \textbf{1.2 Glossário de Termos} — Vocabulário fechado e controlado, sem ambiguidade.
	\end{itemize}
	
	Essas mudanças tornam o documento:
	
	\begin{itemize}[leftmargin=2em]
		\item mais legível;
		\item mais acadêmico;
		\item mais defensável em banca ou revisão técnica.
	\end{itemize}
	
	\subsection{Histórico de Versões Padronizado}
	
	A tabela de histórico foi ajustada para um formato estável e legível, com coluna de \emph{Status}:
	
	\begin{itemize}[leftmargin=2em]
		\item diferencia claramente versões \textbf{congeladas} de versões \textbf{ativas};
		\item facilita auditoria e evolução futura;
		\item alinha-se com boas práticas de RFCs e especificações técnicas.
	\end{itemize}
	
	Embora seja um ajuste pequeno, ele tem \textbf{impacto significativo} em projetos sérios.
	
	\section{Estado Atual do Projeto}
	
	Após os ajustes, o estado real do projeto é o seguinte:
	
	\begin{longtable}{|l|l|}
		\hline
		\textbf{Elemento} & \textbf{Status} \\
		\hline
		Especificação & Completa, consistente e governada \\
		Vocabulário & Fechado \\
		Versionamento & Formal \\
		Processo decisório (ADR) & Definido \\
		Critério de conformidade & Objetivo \\
		Código & Ainda não iniciado (corretamente) \\
		\hline
	\end{longtable}
	
	\bigskip
	\noindent
	\textbf{Conclusão:} o projeto está no \textbf{ponto exato} onde deve estar antes de qualquer linha de código.
	
	\section{Próximo Passo Correto (Engenharia Disciplinada)}
	
	Neste momento, não existem mais dúvidas conceituais.  
	O próximo passo \textbf{não é arbitrário}; ele segue a hierarquia arquitetural já definida.
	
	\subsection{Ordem Correta e Recomendada}
	
	\paragraph{Passo 1 — Criar o ADR001 (Obrigatório)}
	\begin{itemize}[leftmargin=2em]
		\item exercita a governança;
		\item valida o processo decisório;
		\item fixa a primeira decisão derivada diretamente da SPEC.
	\end{itemize}
	
	\paragraph{Passo 2 — Template de Pull Request (Opcional)}
	\begin{itemize}[leftmargin=2em]
		\item baixo custo;
		\item alto impacto futuro.
	\end{itemize}
	
	\paragraph{Passo 3 — Testes de Invariantes (Antes do Código)}
	\begin{itemize}[leftmargin=2em]
		\item garante que a implementação não viole a SPEC;
		\item transforma arquitetura em propriedade verificável.
	\end{itemize}
	
	\paragraph{Passo 4 — Implementação do Nodo-Transistor}
	\begin{itemize}[leftmargin=2em]
		\item primeira unidade concreta;
		\item isolada;
		\item completamente testável.
	\end{itemize}
	
	\section{Recomendação Final}
	
	\begin{quote}
		\textbf{Vamos criar agora o ADR001 completo e oficial.}
	\end{quote}
	
	Esse é o passo que transforma a SPEC em um \textbf{sistema vivo}, e não apenas em um documento estático.
	
	Após a confirmação, serão gerados:
	
	\begin{itemize}[leftmargin=2em]
		\item o arquivo \texttt{ADR001-vector-16-float32};
		\item no padrão definitivo;
		\item pronto para a pasta \texttt{ADRs/};
		\item com status inicial adequado (\emph{Proposta} ou \emph{Aceito}).
	\end{itemize}
	
	\noindent
	A partir daí, o projeto segue exatamente no caminho arquitetural traçado.
	
\end{document}
